\documentclass[10pt]{article}
\usepackage[T1]{fontenc}
\usepackage[utf8]{inputenc}
\usepackage[french]{babel}
\usepackage{amsmath,amssymb,mathrsfs,mathpazo}
\usepackage{geometry}
\geometry{a4paper,margin=2cm}
\usepackage{graphicx,array,tabularx,tabularray,colortbl,multirow,multicol,booktabs,enumitem,pifont,ulem,fancybox,diagbox,cancel}
\usepackage{pstricks,pstricks-add,pst-plot,pst-tree,pst-node,pst-text,pst-func,pst-math,pst-bspline,pst-3dplot,pst-eucl,pst-solides3d}
\let\pscalculate\undefined
\usepackage{pst-bezier}
\usepackage{tikz,pgfplots,tkz-tab}
\usetikzlibrary{arrows,arrows.meta,patterns,calc,positioning,trees,decorations.markings,shapes.geometric,shapes.misc,shapes.arrows,intersections,decorations.pathmorphing,decorations.pathreplacing,backgrounds,shadows,babel}
\pgfplotsset{compat=1.18}
\usepgfplotslibrary{fillbetween}
\usepackage[np]{numprint}
\usepackage{listings,makecell,pas-tableur,twoopt}
\let\barre\undefined
\usepackage[pstricks]{tablvar}
\definecolor{darkgreen}{RGB}{0,100,0}
\providecommand{\red}{\color{red}}
\providecommand{\blue}{\color{blue}}
\providecommand{\green}{\color{green}}
\providecommand{\black}{\color{black}}
\providecommand{\white}{\color{white}}
\providecommand{\gray}{\color{gray}}
\providecommand{\cyan}{\color{cyan}}
\providecommand{\magenta}{\color{magenta}}
\providecommand{\yellow}{\color{yellow}}
\providecommand{\violet}{\color{violet}}
\providecommand{\orange}{\color{orange}}
\providecommand{\euro}{\texteuro}
\providecommand{\EUR}{\texteuro}
\providecommand{\ieme}{\textsuperscript{e}}
\providecommand{\degres}{\ensuremath{^\circ}}
\providecommand{\degree}{\ensuremath{^\circ}}
\providecommand{\vv}[1]{\overrightarrow{#1}}
\providecommand{\Coord}[1]{#1}
\providecommand{\decosix}{}
\providecommand{\starredbullet}{\ensuremath{\bullet}}
\pagestyle{empty}
\setlength{\parindent}{0pt}
\usepackage{adjustbox}\DeclareUnicodeCharacter{2236}{:}\DeclareUnicodeCharacter{2190}{\ensuremath{\leftarrow}}
\providecommand{\R}{\mathbb{R}}
\providecommand{\N}{\mathbb{N}}
\providecommand{\Z}{\mathbb{Z}}
\providecommand{\Q}{\mathbb{Q}}
\providecommand{\D}{\mathbb{D}}
\providecommand{\C}{\mathbb{C}}
\providecommand{\e}{{\mathrm{e}}}
\providecommand{\vect}[1]{\overrightarrow{#1}}
\providecommand{\vectt}[1]{\overrightarrow{\mathrm{#1}}}
\providecommand{\vvt}[1]{\overrightarrow{\mathrm{#1}}}
\providecommand{\np}[2][]{{#2~#1}}
\providecommand{\vv}[1]{{\overrightarrow{#1}}}
\renewcommand{\d}{{\mathrm{d}}}
\providecommand{\degres}{{^\circ}}
\providecommand{\degre}{{^\circ}}
\providecommand{\blue}{\color{blue}}
\providecommand{\red}{\color{red}}
\providecommand{\enskip}{\quad}
\providecommand{\llbracket}{\lbrack\!\lbrack}
\providecommand{\rrbracket}{\rbrack\!\rbrack}
\providecommand{\euro}{{\text{€}}}
\providecommand{\no}{{\text{n°}}}
\providecommand{\interval}[2]{{[#1;#2]}}
\providecommand{\intervFO}[2]{{[#1;#2[}}
\providecommand{\dotfill}{\ldots}

\providecommand{\ds}{\displaystyle}
\providecommand{\textding}[1]{#1}
\providecommand{\Oij}{\(\mathrm{O};\vec{\imath},\vec{\jmath}\)}
\providecommand{\Oijk}{\(\mathrm{O};\vec{\imath},\vec{\jmath},\vec{k}\)}
\providecommand{\Ouv}{\(\mathrm{O};\vec{u},\vec{v}\)}
\providecommand{\intervFF}[2]{\ensuremath{[#1;#2]}}
\renewcommand{\euro}{\ifmmode\text{\texteuro}\else\texteuro\fi}
\DeclareUnicodeCharacter{00B0}{\ifmmode^{\circ}\else\textdegree\fi}
\DeclareUnicodeCharacter{20AC}{\ifmmode\text{\texteuro}\else\texteuro\fi}
% Énoncés : Éducation nationale. Transcription et corrigés : APMEP, Denis Vergès et contributeurs. Découpage et classement : CoopMaths, Sébastien Lozano, Éric Elter et contributeurs. Adaptation de la présentation : Maths Explorer.
% CC BY-SA : https://coopmaths.fr/www/about/
% https://coopmaths.fr/alea/static/e3c/2020/tex/e3c_2020_00_sujet42_2.tex
\begin{document}
\noindent\textit{APMEP / CoopMaths. CC BY-SA. Présentation adaptée par Maths Explorer.}\par\medskip
\medskip

Un fromager fait l'inventaire des produits qu'il a en cave. 

Le graphique ci-dessous indique la répartition de ses trois types de fromages : au lait de chèvre, au lait de vache ou au lait de brebis.

\begin{center}
\psset{unit=1cm}
\begin{pspicture}(-2,-2)(2,3)
\rput(0,2.8){\large Répartition des types de fromage}
\pscircle(0,0){2}
\pswedge(0,0){2}{0}{54}\rput(1.6;35){\footnotesize fromage au }\rput(1.5;20){\footnotesize lait de brebis}
\pswedge(0,0){2}{54}{144}\rput(1.6;90){\footnotesize fromage au lait}\rput(1.3;85){\footnotesize de chèvre}
\pswedge(0,0){2}{144}{360}\rput(1.3;270){fromage au lait}\rput(1.6;270){de vache}
\end{pspicture}
\end{center}

Chacun de ses trois types de fromages se partage en deux catégories: frais ou affiné. Le tableau suivant donne la répartition des fromages de chaque catégorie suivant leur affinage:

\begin{center}
\begin{tabularx}{\linewidth}{|m{3cm}|*{2}{>{\centering \arraybackslash}X|}}\hline
&frais &affiné\\ \hline
Lait de vache & 20\,\% &80\,\% \\ \hline
Lait de chèvre&40\,\%& 60\,\%\\ \hline
Lait de brebis&70\,\% &30\,\%\\ \hline
\end{tabularx}
\end{center}

Le fromager prend un fromage au hasard. On note les évènements suivants :

\setlength\parindent{1cm}
\begin{itemize}[label=\textbullet]
\item $V$ : \og le fromage est fait avec du lait de vache \fg{};
\item $C$ : \og le fromage est fait avec du lait de chèvre \fg{} ;
\item $B$ : \og le fromage est fait avec du lait de brebis \fg{} ;
\item $F$ : \og le fromage est frais \fg{} ;
\item $A$ : \og le fromage est affiné \fg.
\end{itemize}
\setlength\parindent{0cm}

\medskip

\begin{enumerate}
\item Donner les probabilités $P_C(A)$ et $P(B)$.
\item Démontrer que $P(A) = 0,675$.
\item Le fromager prend au hasard un fromage affiné. Quelle est la probabilité qu'il
s'agisse d'un fromage au lait de vache? On donnera le résultat à $10^{-3}$ près.
\end{enumerate}

\bigskip

\end{document}
